% JMLR MLOSS manuscript for pyvinecopulib.
%
% A skeleton: each section holds a plan box -- the argument, the evidence, the
% displays and what the section leaves out -- where its prose will go.
% Table 1 and Listing 1 are final; Figure 1 and Table 2 are placeholders of
% their target size, so the PDF already shows where the floats put page 4.
% plan.md, beside this file, holds the reasoning. Build with ./build.sh.
\documentclass[twoside,11pt]{article}

\usepackage{jmlr2e}

% Packages are loaded after jmlr2e, which already loads amssymb, natbib and
% graphicx.
\usepackage{booktabs}
\usepackage{listings}
\usepackage{tcolorbox}
\usepackage{lastpage}

% A plan box stands in for a section's prose until it is written.
\newtcolorbox{plan}[1]{%
  colback=black!3, colframe=black!40, coltitle=white,
  fonttitle=\bfseries\small, fontupper=\small,
  left=4pt, right=4pt, top=2pt, bottom=2pt,
  title={Plan: #1}}

\lstset{%
  language=Python, basicstyle=\ttfamily\footnotesize, columns=fullflexible,
  keepspaces=true, upquote=true, numbers=left, numberstyle=\tiny,
  xleftmargin=1.5em, frame=tb, captionpos=b}

\newcommand{\pkg}[1]{\textsf{#1}}

% {volume}{year}{pages}{date submitted}{date published}{paper id}{authors};
% JMLR fills in the rest on acceptance.
\jmlrheading{}{2026}{1-\pageref{LastPage}}{}{}{}{Thibault Vatter and Thomas
Nagler}

\ShortHeadings{Composable Vine Copula Distributions}{Vatter and Nagler}
\firstpageno{1}

\begin{document}

\title{pyvinecopulib: Composable Vine Copula Distributions for Scientific
Python and Machine Learning}

\author{\name Thibault Vatter \email thibault.vatter@hesge.ch \\
       \addr Geneva School of Business Administration\\
       HES-SO University of Applied Sciences and Arts Western Switzerland\\
       1227 Carouge, Switzerland
       \AND
       \name Thomas Nagler \email thomas.nagler@stat.uni-muenchen.de \\
       \addr Department of Statistics\\
       Ludwig-Maximilians-Universit\"at M\"unchen\\
       80539 M\"unchen, Germany}

\editor{TBD}

\maketitle

\begin{abstract}%
Vine copulas build flexible multivariate distributions from bivariate
building blocks, and the C++ library vinecopulib provides mature estimation
for them. \pkg{pyvinecopulib} turns these models into components of Python
machine-learning workflows. Its core is four array-agnostic contracts---for
pair copulas, vines, univariate margins and joint distributions on the data
scale---with base classes that implement structure selection, the evaluation
and simulation cascades and two-step estimation once, against those
contracts. Compiled vinecopulib objects, NumPy and PyTorch implementations and
user-defined parts all satisfy the same contracts, so each inherits the same
algorithms, while the pair copulas hosted in a vine decide how each tree level
executes: as compiled calls, on a thread pool, or as one stacked GPU
operation. On this foundation the library provides scikit-learn estimators and
differentiable PyTorch modules, validated against the compiled engine.
\end{abstract}

\begin{keywords}
  vine copulas, density estimation, PyTorch, scikit-learn, open-source
  software
\end{keywords}

\section{Introduction}

\begin{plan}{about 0.7 page, plus Table~\ref{tab:related}}
\textbf{Argument.} Vine copulas \citep{bedford2002vines,aas2009pair} are
flexible density models used in machine learning for tabular density
estimation, synthetic data, conditional simulation and learned dependence.
Python offers narrow vine packages and differentiable libraries without vines
(Table~\ref{tab:related}). A binding to the C++ engine alone cannot host
user-defined or learnable parts, run on accelerators, or sit in scikit-learn
and PyTorch pipelines. State the central phrase once: \emph{the vine
algorithms are written once, against four contracts; every implementation of
a part---compiled, NumPy, PyTorch or the user's own---inherits them, and
decides how its tree levels execute.} Then the contribution paragraph and
three bullets (plan.md, section 2.4).

\textbf{Evidence.} Table~\ref{tab:related}; dependents named in the cover
letter.

\textbf{Citations.} Methods, engine and R interface:
\citet{nagler2026rvinecopulib}; statistical algorithms:
\citet{nagler2026algorithms}; the earlier computational-graph formulation of
the PyTorch lane: \citet{cheng2025vine}.

\textbf{Left out.} Vine theory beyond three sentences; pair-copula families;
selection algorithms.
\end{plan}

\begin{figure}[t]
\centering
\fbox{\begin{minipage}[c][0.3\textheight][c]{0.95\textwidth}
\centering\footnotesize
\textbf{Placeholder for Figure~\ref{fig:architecture}: full width, about 0.4
page with its caption.}\\[6pt]
\begin{minipage}{0.92\textwidth}
\raggedright
\textbf{Consumers:} scikit-learn (\texttt{VineDensity},
\texttt{VineRegressor}); PyTorch (\texttt{nn.Module} composition,
\texttt{.to(device)}, autograd, optimizers); NumPy/SciPy code.\\[3pt]
\textbf{Waist, the four contracts (stable since 1.0):} \texttt{MarginLike}
and \texttt{VinecopLike} feed \texttt{VinedistLike}; \texttt{BicopLike} is
hosted at each (tree, edge) of \texttt{VinecopLike}.\\[3pt]
\textbf{Implementations, one column each:} compiled from vinecopulib
(\texttt{Bicop}, \texttt{Vinecop}, \texttt{Kde1d}); generic bases
(\texttt{BicopBase}, \texttt{VinecopBase}, \texttt{MarginBase},
\texttt{VinedistBase} $\to$ \texttt{Vinedist}); PyTorch lane
(\texttt{TorchTllBicop}, \texttt{TorchVinecop}, \texttt{TorchKde1d},
\texttt{TorchVinedist}); adapters and user parts (dashed).\\[3pt]
\textbf{Inset beside \texttt{VinecopBase}:} tree levels, and dependency waves
for the inverse; a level runs as one stacked call when its pairs' class
supplies one, otherwise pair by pair, and is fitted on a thread pool when the
pair releases the GIL.\\[3pt]
\textbf{Foundation:} vinecopulib, kde1d, wdm (C++17, Eigen), shared with
rvinecopulib.
\end{minipage}
\end{minipage}}
\caption{Architecture of \pkg{pyvinecopulib}. Consumers depend only on four
contracts; each contract has compiled, generic, PyTorch and user-defined
implementations, and one vine cascade runs over all of them, a tree level at a
time.}
\label{fig:architecture}
\end{figure}

\begin{table}[t]
\centering
\footnotesize
\setlength{\tabcolsep}{4pt}
\begin{tabular}{@{}lccccc@{}}
\toprule
 & \pkg{pyvinecopulib} & \pkg{VineCopulas} & \pkg{torchvinecopulib}
 & \pkg{mixedvines} & \pkg{Copulas} \\
 & 1.1 & 2.0.3 & 1.1.2 & 1.3.5 & 0.14.1 \\
\midrule
R-vines with selection & \checkmark & \checkmark & \checkmark & C-vine
 & partial\textsuperscript{a} \\
Pair copulas (param.\ / nonpar.) & 11 / TLL & 6 / -- & -- / TLL, KDE
 & 3 / -- & 3 / -- \\
Discrete and mixed data & exact & jittered & -- & exact\textsuperscript{b}
 & -- \\
Joint density and sampling & \checkmark & -- & \checkmark & \checkmark
 & sampling \\
scikit-learn estimators & \checkmark & -- & -- & -- & -- \\
Autograd and GPU & \checkmark & -- & \checkmark & -- & -- \\
User pair copulas in a vine & \checkmark\textsuperscript{c} & -- & --
 & \checkmark & -- \\
License & MIT & GPL-3.0 & MIT & GPL-3.0 & BUSL-1.1 \\
\bottomrule
\end{tabular}
\caption{Python packages implementing vine copulas, as of October 2026
\citep{claassen2024vinecopulas,torchvinecopulib,onken2016mixed,copulas}.
\textsuperscript{a}Three trees by default; its constructor warns that it is
not fully tested. \textsuperscript{b}C-vines only. \textsuperscript{c}Including
learnable and non-simplified pair copulas. Statistical capabilities shared
with the R interface are compared in \citet{nagler2026rvinecopulib}.}
\label{tab:related}
\end{table}

\section{Design}
\label{sec:design}

\begin{plan}{about 1.15 pages in total, one subsection per box of
Figure~\ref{fig:architecture}}
\textbf{Left out.} The internals of the engine's computation index; the wave
algorithm; the kernels. These belong to \citet{nagler2026rvinecopulib} and to
the methods paper.
\end{plan}

\subsection{Four contracts}

\begin{plan}{about 0.4 page}
\textbf{Argument.} Each contract is a protocol plus a base. The compiled
classes satisfy the protocols structurally, because a nanobind class cannot
inherit a Python ABC; the bases derive everything else from a few primitives;
a member a part cannot serve has a default that raises by name. Fitting is
declared rather than hooked: \texttt{bicop\_class}, \texttt{vinecop\_class},
\texttt{margin\_class} and \texttt{controls\_class} name the parts, so
\texttt{from\_data} needs no callback.

\textbf{Evidence.} The protocol and contract test suites; the import-surface
test that pins the layers.
\end{plan}

\subsection{Algorithms written once}

\begin{plan}{about 0.4 page}
\textbf{Argument.} \texttt{VinecopBase} implements the evaluation and
simulation cascades, conditional sampling and structure selection, which
reproduces the engine's selected matrix \citep{dissmann2013selecting};
\texttt{VinedistBase} implements two-step estimation and serialization.
Discrete edges reach a pair through its declared variable types, conditional
vines through a context. Arrays are placed by inference, so a part written for
NumPy or PyTorch, on any device and dtype, needs no conversion code.

\textbf{Evidence.} A vine hosting \texttt{Bicop} selects the same structure
as \texttt{Vinecop} and agrees on densities to $10^{-12}$
(Table~\ref{tab:fidelity}).
\end{plan}

\subsection{Execution by tree level}

\begin{plan}{about 0.35 page}
\textbf{Argument.} The engine's computation index
\citep[Section~3.1]{nagler2026rvinecopulib} wires whole tree levels; the
inverse Rosenblatt transform runs in dependency waves. A level runs as one
stacked call when its pairs' class supplies one, otherwise pair by pair, with
results identical to edge-by-edge evaluation. A fit is stacked, threaded, or
edge by edge, as the pair class allows. One sentence: an earlier version of the
PyTorch lane followed the computational-graph formulation of
\citet{cheng2025vine}, and execution by level replaced it for efficiency.

\textbf{Evidence.} Table~\ref{tab:fidelity}: the stacked and pair-by-pair rows, and the threaded fit in its caption.
\end{plan}

\section{Vine Distributions in Machine-Learning Workflows}

\begin{plan}{about 0.55 page, including Listing~\ref{lst:learnable}}
\textbf{Argument.} scikit-learn \citep{pedregosa2011scikit}:
\texttt{VineDensity} and \texttt{VineRegressor} take
\texttt{distribution=} a class, accept mixed DataFrames, and pass
\texttt{check\_estimator}. PyTorch \citep{paszke2019pytorch}: every part is an
\texttt{nn.Module}, so \texttt{.to}, \texttt{state\_dict} and autograd reach
the whole distribution, discrete and trainable margins included. Margins from
SciPy, \texttt{torch.distributions} or user code arrive through adapters.
Listing~\ref{lst:learnable} composes a user-defined pair, the stock vine and
fitted kernel margins into one trainable module. Two sentences follow it: a
network of covariates in place of the constant correlation gives a
non-simplified vine, and naming \texttt{bicop\_class} and
\texttt{vinecop\_class} on subclasses lets \texttt{from\_data} and
\texttt{VineDensity} fit such a model.

\textbf{Left out.} API enumeration; controls.
\end{plan}

\begin{lstlisting}[float=t, caption={A learnable pair copula, hosted by the
stock vine and composed with fitted kernel margins into a data-scale PyTorch
module (\texttt{data} is an $n \times d$ array; the pair's formulas are in
Appendix~\ref{app:code}).}, label={lst:learnable}]
import torch
from pyvinecopulib.core import BicopBase
from pyvinecopulib.torch import TorchVinecop, TorchVinedist


class Gaussian(BicopBase[torch.Tensor], torch.nn.Module):  # learnable pair
  def _pdf_raw(self, u): ...  # density and the two conditional
  def _hfunc1_raw(self, u): ...  # distribution functions; parameter
  def _hfunc2_raw(self, u): ...  # and constructor in Appendix C


y = torch.as_tensor(data, device="cuda")  # raw data, on the GPU
fit = TorchVinedist.from_data(y)  # kernel margins + vine
pairs = [[Gaussian() for _ in range(d - 1 - t)] for t in range(d - 1)]
copula = TorchVinecop(pairs, fit.vinecop.structure)  # stock vine, user pairs
dist = TorchVinedist(copula, fit.margins)  # an nn.Module
opt = torch.optim.Adam(dist.parameters(), lr=0.05)
for _ in range(200):  # copula MLE, data scale
  opt.zero_grad()
  (-dist.log_prob(y).mean()).backward()
  opt.step()
\end{lstlisting}

\section{Fidelity and Execution}

\begin{plan}{about 0.4 page, including Table~\ref{tab:fidelity}}
\textbf{Argument.} Implementations of the same contract agree, and each
execution strategy wins somewhere: threads at small $n$, the GPU once levels
are wide and $n$ is large. One sentence on tests, coverage, CI on five
platforms, static typing and documentation.

\textbf{Evidence.} Table~\ref{tab:fidelity}, measured on the 1.1.0 wheels in
float64 with paired repetitions.

\textbf{Left out.} Comparisons with R packages; the GPU analysis.
\end{plan}

\begin{table}[t]
\centering
\footnotesize
\setlength{\tabcolsep}{4pt}
\begin{tabular}{@{}lccccc@{}}
\toprule
 & Max.\ rel.\ error & \multicolumn{2}{c}{\texttt{logpdf} (ms)}
 & \texttt{sample} (ms) & \texttt{from\_data} (s) \\
\cmidrule(lr){3-4}
Implementation & vs \texttt{Vinecop} & $n = 10^3$ & $n = 2 \cdot 10^4$
 & $n = 2 \cdot 10^4$ & $n = 2 \cdot 10^4$ \\
\midrule
\texttt{Vinecop}, 1 thread & reference & -- & -- & -- & -- \\
\texttt{Vinecop}, all cores & -- & -- & -- & -- & -- \\
\texttt{VinecopBase} + \texttt{Bicop}, all cores & -- & -- & -- & -- & --
 \\
\texttt{TorchVinecop}, CPU & -- & -- & -- & -- & -- \\
\texttt{TorchVinecop}, CUDA, pair by pair & -- & -- & -- & -- & -- \\
\texttt{TorchVinecop}, CUDA, stacked & -- & -- & -- & -- & -- \\
\quad and compiled & -- & -- & -- & -- & n/a \\
\bottomrule
\end{tabular}
\caption{One \texttt{tll} vine ($d = 10$), seven implementations: agreement
with \texttt{Vinecop} on held-out points, and time. Stacked and pair-by-pair
levels agree to rounding, and their inverse transforms exactly; a threaded fit
equals the serial one. Placeholder: to be measured on the 1.1.0 release.}
\label{tab:fidelity}
\end{table}

\section{Related Software and Availability}

\begin{plan}{about 0.25 page}
\textbf{Argument.} How to read Table~\ref{tab:related}; complementary
libraries without vines (normalizing flows, NumPyro's Gaussian copula,
scikit-learn's kernel density and Gaussian mixture estimators).
\pkg{torchvinecopulib} accompanies \citet{cheng2025vine}, which the authors
co-wrote, and installs \pkg{pyvinecopulib}. Availability: MIT license;
PyPI and conda-forge; documentation at
\url{https://pyvinecopulib.readthedocs.io}; the reviewed version 1.1.0 and its
DOI \citep{pyvinecopulib}; pointers to \citet{nagler2026rvinecopulib} and the
methods paper.

\textbf{Left out.} Roadmap promises.
\end{plan}

\acks{To be written, including the disclosure of funding.}

\newpage

\appendix

\section{Extended Comparison with Python Software}
\label{app:comparison}

\begin{plan}{about 0.75 page}
The full verified matrix behind Table~\ref{tab:related}, with a pinned source
for every cell (plan.md, section 6).
\end{plan}

\section{Fidelity and Timing Setup}
\label{app:setup}

\begin{plan}{about 0.75 page}
Hosts, library versions, repetitions and confidence intervals, and the full
tables behind Table~\ref{tab:fidelity}, including peak memory (plan.md,
section 7).
\end{plan}

\section{Complete Code for Listing~\ref{lst:learnable}}
\label{app:code}

\begin{plan}{about 0.75 page}
The \texttt{Gaussian} pair with its parameter, constructor and formulas, as
executed in notebook 10; and a vine hosting \texttt{Bicop} that reproduces
\texttt{Vinecop}.
\end{plan}

\section{From the 0.x Bindings to 1.x}
\label{app:changes}

\begin{plan}{about 0.5 page}
The improvements and extensions since the bindings, mirroring the
repository's citing page, which MLOSS asks the cover letter to point to.
\end{plan}

\vskip 0.2in
\bibliography{references}

\end{document}
